\section*{ЗАКЛЮЧЕНИЕ}
\addcontentsline{toc}{section}{ЗАКЛЮЧЕНИЕ}

В ходе выполнения данной работы были изучены теоретические основы обработки сигналов в беспроводных системах связи, а также проанализированы методы оценивания канала связи.

В первом разделе были рассмотрены базовые модели распространения радиоволн. Показано, что в свободном пространстве мощность сигнала убывает пропорционально квадрату расстояния, однако влияние подстилающей поверхности (двулучевая модель) приводит к более быстрому затуханию — пропорционально четвертой степени расстояния (формула Введенского). Также были детально изучены крупномасштабные и мелкомасштабные замирания сигнала. Установлено, что многолучевое распространение приводит к временной и угловой дисперсии, а статистические свойства замираний описываются законами Релея и Райса. Знание этих статистических свойств критически важно, так как они определяют вероятность битовой ошибки и предельную пропускную способность системы.

Во втором разделе был рассмотрен практический подход к решению проблемы искажения сигнала. Была сформулирована задача оценивания импульсной характеристики (ИХ) многолучевого канала. Показано, что для компенсации межсимвольных искажений приемное устройство должно знать коэффициенты передачи канала. Для их нахождения используется метод максимального правдоподобия и метод наименьших квадратов. Было установлено, что точность оценивания ИХ напрямую зависит от отношения сигнал/шум, длины обучающей последовательности и степени ортогональности обучающих векторов.

Таким образом, можно сделать вывод, что качественная оценка канала связи является фундаментальной основой для работы алгоритмов эквализации (выравнивания) сигнала. В современных MIMO-системах и сетях 5G, работающих в условиях сложной городской застройки, методы оценивания канала позволяют эффективно бороться с замираниями и обеспечивать высокую скорость передачи данных. Дальнейшим направлением развития данной темы может стать изучение адаптивных алгоритмов оценивания, работающих в реальном времени без использования длинных пилотных сигналов.